\documentclass[10pt,a4paper]{amsart}
\usepackage[margin=19mm]{geometry}
\usepackage{amsmath,amssymb,mathtools}
\usepackage[scr=rsfs,cal=euler]{mathalfa}
\usepackage{xfrac}
\usepackage[unicode,hidelinks,bookmarks=false]{hyperref}
\newcommand{\ie}{i.e.\ }
\newcommand{\cf}{cf.\ }
\newcommand{\resp}{respectively\ }
\newcommand{\Subsection}{\subsection}
\providecommand{\nobreakdash}{\nobreak\hbox{-}}
\setlength{\emergencystretch}{2em}
\multlinegap0pt
\usepackage{algorithm}\usepackage{algpseudocode}
\usepackage{mathtools}
\usepackage{xcolor}
\usepackage[shortlabels]{enumitem}
\usepackage{amssymb}
\newtheorem{thm}{Theorem}
\newtheorem{pf}{Pf}
\newcommand{\KL}[2]{D_{\mathrm{KL}}(#1 \parallel #2)}
\newcommand{\R}{{\mathbb R}}
\newcommand{\Rd}{\R ^{d}}
\newcommand{\rmd}{{\rm d}}
\title{Unbalanced Optimal Transport and Density Control for Discrete-Time Linear Systems}
\author{Haruto Nakashima, Siddhartha Ganguly, Kenji Kashima}
\date{Source: arXiv:2605.06391v1 (CC BY 4.0)}
\begin{document}
\maketitle

\noindent\emph{Source and licence.} Unbalanced Optimal Transport and Density Control for Discrete-Time Linear Systems. Version arXiv:2605.06391v1 (CC BY 4.0), distributed by arXiv under the Creative Commons Attribution 4.0 International licence, http://creativecommons.org/licenses/by/4.0/, as recorded in the arXiv licence metadata for this exact version and shown on the abstract page https://arxiv.org/abs/2605.06391v1. Full licence notice: \texttt{licenses/arxiv-2605.06391-CC-BY-4.0.txt}.\par\smallskip

\noindent\textit{Editorial scope.} Counted unit: the article's thm thm:uot-global (source0.tex lines 631--633). The article's complete original proof of it is delivered with it (source0.tex lines 635--651, one \texttt{proof} environment of 17 lines). No mathematics is written, repaired or re-derived: the statement and the proof are the article's own lines, in the article's own notation, and no retained line is substituted or re-flowed.  Retained uncounted as the source-local context the proof needs: lines 484--491; lines 625--627; lines 657--670.  Nothing was omitted from any retained span, so the package carries no omission record.  No retained line contains a citation, so the wrapper carries no bibliography and no bibliography entry.  No retained line contains a figure environment, an image-inclusion command or drawing code, so no image is delivered and the package declares no asset dependency.  Editorial formatting only, in addition: an AMS \texttt{amsart} preamble that restates the article's own declarations and macros used by the retained lines, namely the environments \texttt{thm}, \texttt{pf} on separate counters, together with the standard packages those lines need.  The retained environments print in this document as Theorem 1, Pf 1 in order of appearance, whereas the article prints the same items with the numbering of its own full document.  The complete native source of the article as distributed in the arXiv e-print for this exact version is bundled unchanged as \texttt{sources/200017/source0.tex}.
\begin{equation}\label{eq:uot}
\begin{aligned}
\inf_{\pi} 
\int_{\Rd}\int_{\Rd} & \|x_2-x_1\|^2 \,\rmd \pi(x_1,x_2) \\
&+ \gamma \KL{\pi_1}{\alpha}
+ \gamma \KL{\pi_2}{\beta}.
\end{aligned}
\end{equation}

\begin{align}\label{eq:uot-subproblem}
    \min_{m_1,\Sigma_1,m_2,\Sigma_2}\quad M(m_1,m_2) + C(\Sigma_1,\Sigma_2).
\end{align}

\begin{algorithm}[h]
\caption{UOT between Gaussians}
\label{alg:uot}
\begin{enumerate}
\item \textbf{Mean and covariance optimization (convex).}
Solve~\eqref{eq:uot-subproblem} \\
\(p^*\longleftarrow\) Optimal value of~\eqref{eq:uot-subproblem}
\item \textbf{Mass computation (closed form).}
\begin{align}\label{eq:c-star}
    c^\ast \longleftarrow &\sqrt{c_\alpha c_\beta} \exp\left(-\frac{p^\ast}{2\gamma} - \frac{L}{4} \right), \\
    &(L := \log\det(\Sigma_\alpha)+\log\det(\Sigma_\beta)-2d.) \nonumber
\end{align}
\end{enumerate}
\end{algorithm}

\begin{thm}[Global optimality of Algorithm~\ref{alg:uot}]\label{thm:uot-global}
Algorithm~\ref{alg:uot} yields a globally optimal solution to the UOT problem~\eqref{eq:uot}.
\end{thm}

\begin{pf}
Let $p^\ast$ denote the optimal value of~\eqref{eq:uot-subproblem}. Since
\begin{align}
&c\,M(m_1,m_2) + c\,C(\Sigma_1,\Sigma_2) + \psi(c) \\
= &c\big(M(m_1,m_2)+C(\Sigma_1,\Sigma_2)\big) + \psi(c),
\end{align}
we have
\begin{align}
\inf_{c,m_1,\Sigma_1,m_2,\Sigma_2}
\big\{c\,M + c\,C + \psi(c)\big\}
=
\inf_{c>0}\ \Big\{ c\,p^\ast + \psi(c)\Big\}.
\label{eq:separation}
\end{align}
The right-hand side is a one-dimensional convex optimization in $c$; hence the first-order optimality condition
gives the unique minimizer in closed form, which is exactly~\eqref{eq:c-star}. \qed
\end{pf}

\end{document}
